% =====================================================================
%  SUGGESTED ADDITIONS / REPLACEMENTS FOR sec:rt-partition
%  Incorporating Toro (2026), bioRxiv 10.64898/2026.05.14.725207
%
%  Each block is marked [ADD], [REPLACE] or [EDIT] with the target
%  location in RT0-RT7_theory.txt.
% =====================================================================


% ---------------------------------------------------------------------
% [ADD] Preamble: new provenance macro
% Place alongside \derived, \inherit, \noprov
% ---------------------------------------------------------------------
\newcommand{\operational}{\textsc{operational}}

% ---------------------------------------------------------------------
% [ADD] Preamble: \todo for the verification markers used below.
% If you already load todonotes, delete this block -- \providecommand
% will not override it, but the margin style may differ.
%   \usepackage[colorinlistoftodos,textsize=footnotesize]{todonotes}
% Minimal fallback (no package required):
% ---------------------------------------------------------------------
\providecommand{\todo}[1]{%
  \marginpar{\raggedright\footnotesize\textbf{TODO:}~#1}}

% Every \todo{} inserted below marks a quantity RECOMPUTED BY ANALYSIS
% rather than read directly from the source. See the VERIFICATION
% REGISTER at the end of this file. Strip all \todo{} before submission.


% ---------------------------------------------------------------------
% [ADD] Bibliography entry
% ---------------------------------------------------------------------
% @article{toro2026,
%   author  = {Toro, Nicol\'{a}s},
%   title   = {Landscape of retron diversity across the {SPIRE} prokaryotic
%              metagenome resource reveals candidate novel type {XI}-like lineages},
%   journal = {bioRxiv},
%   year    = {2026},
%   doi     = {10.64898/2026.05.14.725207},
%   note    = {Preprint, posted 14 May 2026}
% }


% =====================================================================
% [ADD] New row for Table~\ref{tab:rt-sources}
% Insert after the mestre2020 row, before arkhipova2023
% =====================================================================

Toro~\cite{toro2026}
  & 107{,}078 SPIRE representative genomes $\rightarrow$ 1{,}286 candidates;
    1{,}928 reference RTs~\cite{mestre2020}
  & Profile HMM, phylogenetic placement
  & First large-scale metagenomic retron survey. \emph{Abandons RT0}: the span
    is RT1--RT7 throughout, with no statement of the change and no mention of
    RT0. The span is realised as the 203 alignment columns surviving a
    ${>}50\%$ gap filter. Detection uses a general HMM built on the full
    reference sequences, so two coordinate systems operate in one study
  & \operational\tnote{f} \\[2pt]

% [ADD] corresponding tablenote, after note [e]
% \item[f] A criterion is published, but it is not a criterion about the
%   partition: a generic gap-occupancy threshold produces a block of columns
%   which is then named after RT1--RT7. No source is cited for the RT1--RT7
%   span, and no reason is given for the departure from RT0--RT7.


% =====================================================================
% [REPLACE] Final paragraph of \subsection{Origin and propagation of the
% partition}, from ``Toro \textit{et al.}~\cite{toro2019} and Rodr{\'i}guez
% Mestre...'' through ``...applied to the boundary question.''
% =====================================================================

Toro \textit{et al.}~\cite{toro2019} and Rodr{\'i}guez Mestre \textit{et
al.}~\cite{mestre2020} then use ``RT0--7'' as a given. Both name individual
elements when locating the retron-specific regions---region~X between RT2 and
RT3, region~Y within RT7---and neither refers to RT1 anywhere in its published
main text.

Two departures from this pattern of inheritance have since occurred, and they
are of different kinds. The first is the myRT reference package~\cite{sharifi2022},
whose profile models are built on the Pfam RVT\_1 domain (PF00078) rather than
on the RT0--RT7 span; a coordinate system independent of the partition has
therefore been in use since 2022, though it has not been applied to the
boundary question. The second is more consequential, because it comes from
within the retron lineage itself. Toro~\cite{toro2026} surveys 107{,}078 SPIRE
representative genomes using the 1{,}928-sequence reference set
of~\cite{mestre2020}, and throughout refers to the \emph{RT1--RT7} motif core:
as the basis of the ``RT domain completeness'' criterion used to assign
confidence tiers, and as the region retained for the maximum-likelihood
reference tree. RT0 does not appear in the paper. No reason is given for its
removal, and no source is cited for the RT1--RT7 span that replaces it. The
element whose provenance in the retron literature this chapter traces has
been dropped by the lineage that carried it, without comment.

The manner of the replacement matters as much as the fact of it. The span is
not delimited by inspection, as in~\cite{xiong1990}, nor adopted silently, as
in~\cite{toro2014,toro2019,mestre2020}. It is stated as ``the 203 positions
corresponding to the conserved RT1--RT7 motif core'', obtained ``by removing
columns containing more than 50\% gaps''~\cite{toro2026}. That sentence admits
two readings. Under the first, the gap filter alone produces the 203 columns,
which are then named after RT1--RT7---in which case the partition is a label
applied to the output of a generic alignment filter, and the boundary is
wherever gap occupancy happens to fall for the particular set of sequences
aligned. Under the second, the region was first delimited as RT1--RT7 by
unstated means and the gap filter applied within it, in which case the
boundary is again inherited and again without criterion. The two readings
imply different objects, and neither supplies a criterion specific to the
partition. This is recorded in Table~\ref{tab:rt-sources} as a distinct
provenance category, \operational, rather than as inheritance: a criterion is
published; it is simply not a criterion about the elements it names.

The consequence is a coordinate system that is now internally plural. In
\cite{toro2026}, detection is performed with a general profile HMM built from
the full 1{,}928 reference sequences, while phylogenetic inference is
performed on the 203-column block. Two spans are therefore in use in a single
study for two different tasks, without the relationship between them being
stated. Section~\ref{sec:rt-partition-elements} argues that this is the right
instinct---the elements are not equally transferable and different analyses
have different requirements---but it is exercised here implicitly rather than
as a stated design choice, and the mapping between the two systems is not
published.


% =====================================================================
% [REPLACE] The paragraph beginning ``A related inconsistency affects any
% attempt to reproduce these datasets.'' --- keep it, and append the
% following continuation.
% =====================================================================

The threshold has since changed form again, and this time its behaviour is
reported. Toro~\cite{toro2026} replaces the length criterion with a bit-score
criterion: hits from a general retron RT HMM are filtered at an E-value of
$1\times10^{-10}$, yielding 31{,}280 initial hits, and sequences scoring at
least 160 bits are retained as candidates ($n = 1{,}319$). The cutoff is
calibrated against the reference set itself, and the calibration is published:
the lowest-scoring reference sequence reaches 114.5 bits, and 195 of the
1{,}928 reference sequences ($10.1\%$) fall below the 160-bit cutoff,
predominantly from type~VI ($n = 91$), type~I-B2 ($n = 36$), type~I-B1 ($n =
34$) and type~IV ($n = 20$). The author states that the threshold is a
conservative trade-off and will underrepresent the most divergent lineages.
This is a material improvement on the position described above: a filter is
applied to a stated object, and the fraction of the reference set it excludes
is quantified rather than left to be inferred. It should be read as such. What
it does not do is bear on the boundary question. A score threshold on a
profile built from whole reference sequences says nothing about where the
elements of the partition begin and end, and the 1{,}928-sequence reference
set to which it is applied was itself assembled under the ${\ge}200$-amino-acid
filter applied to the RT0--7 span by~\cite{toro2019,mestre2020}. That filter
therefore remains upstream of every result in~\cite{toro2026}, defined on a
span the analysis no longer uses.


% =====================================================================
% [REPLACE] The RT0 paragraph in \subsection{Element-by-element}, from
% ``RT0 is the element the retron literature reports in and cannot align.''
% Replace the final sentence and append.
% =====================================================================

% ... keep the paragraph as written up to and including:
% ``...therefore measures something whose starting point is undefined for a
% quarter of the dataset.''
% then append:

The span has since been redrawn to exclude it. Toro~\cite{toro2026} works
throughout with an RT1--RT7 core and does not mention RT0 at any point. Read
against the evidence assembled above---absent from the original
partition~\cite{xiong1990}, added from a different retroelement
literature~\cite{malik1999,zimmerly2001}, delimited in LtrA only as the
content of a proteolytic fragment~\cite{blocker2005}, excluded from every
published transferability count~\cite{simon2008}, and not alignable with the
group~II reference in any retron representative of that
study~\cite{simon2008}---the removal is unsurprising. What is notable is that
it is unannounced. An element that entered the retron framework without a
stated source has left it without a stated reason, and in between it
determined which sequences passed the length filter that defines the reference
set still in use. Neither its entry nor its exit was accompanied by a claim
that could be checked, which is the pattern this chapter documents rather than
an isolated lapse.


% =====================================================================
% [EDIT] RT1 discussion, \subsection{Element-by-element}
% Amend the sentence recording that RT1 is unmentioned in the retron literature
% =====================================================================

% Current: ``...no reference to RT1 in \cite{toro2019,mestre2020} (inspection
% of the published text)''
% Suggested replacement:

no reference to RT1 in~\cite{toro2019,mestre2020} (inspection of the published
main text); named in~\cite{toro2026} as the opening element of the RT1--RT7
core, without a stated start position or a cited source for the span


% =====================================================================
% [EDIT] Table~\ref{tab:rt-elements}: amended rows
% =====================================================================

RT0 / NTE & Extended fingers
  & Absent from~\cite{xiong1990}; described in non-LTR RTs~\cite{malik1999,zimmerly2001};
    within a 10\,kDa N-terminal proteolytic fragment of LtrA (M1--R85),
    conserved Ala39~\cite{blocker2005}; \emph{not alignable with the group~II
    reference in any of the three retron or three DGR
    representatives}~\cite{simon2008}; excluded from all published
    transferability counts~\cite{simon2008}; \emph{excluded from the span
    altogether in}~\cite{toro2026}, without comment
  & Loop within an N-terminal extension~\cite{yang2025}. Helical character
    predicted in~\cite{blocker2005} \\[3pt]

% ---------------------------------------------------------------------

Region X$^{b,c}$ & Between RT2 and RT3 &
  \motif{NAXXH}, identified in 59\% of 102 retron/retron-like sequences; His
  100\%, Ala 87\%, Asn 62\% conserved~\cite{toro2014}; suggested to be required
  for initiation or catalysis~\cite{mestre2020}; \emph{used as a per-sequence
  completeness filter} in~\cite{toro2026}, where 184 of 207 protease-less
  type~XI-like candidates (88.9\%) are reported to retain it, still with no
  matching rule stated &
  Not resolved \\[3pt]

Region Y$^{b,c}$ & Within RT7 &
  \motif{VTG}, Gly most conserved; 47\% of the same set, with a further 26\%
  carrying an I/L variant at the first position~\cite{toro2014}; implicated in
  targeting the RT to its cognate \textit{msr}~\cite{inouye1999,mestre2020};
  applied jointly with region~X as a filter in~\cite{toro2026} &
  Not resolved \\


% =====================================================================
% [ADD] Optional short paragraph for the close of
% \subsection{Element-by-element}, or for the chapter conclusion.
% This is the observation that generalises the chapter.
% =====================================================================

One asymmetry in current practice is worth recording, because it shows that
the distinction this chapter needs is already available to the field.
Toro~\cite{toro2026} states explicitly that retron types lacking
experimentally validated \textit{msr}/\textit{msd} structures were not
represented in the covariance-model panel, and that the absence of a hit in
those groups should not be read as evidence for the absence of an associated
ncRNA. That is precisely the distinction Simon and Zimmerly~\cite{simon2008}
draw for the protein elements when they note that sufficient residues are
present to correspond to the domains that cannot be aligned. The caution is
applied to the RNA and not to the partition: an element that cannot be aligned
is dropped from the span, while an RNA that cannot be detected is flagged as
undetected. Nothing in the evidence supports treating the two cases
differently.


% =====================================================================
% [ADD] Reference figures to quote when citing \cite{toro2026}
% =====================================================================
% Panel of 28 type-specific HMMs; per-group thresholds set at
%   (pos_min + neg_max)/2, ranging 177.2 bits (type I-B2) to 317.5 bits (clade11-A).
% Reference tree: MAFFT v7; 203 columns; IQ-TREE 2 v2.3.2; ModelFinder selected
%   Q.pfam+F+I+G4; 1,000 ultrafast bootstrap replicates.
%   (cf. \cite{mestre2020}: MAFFT + FastTree, WAG, discrete gamma, 20 categories.)
% Placement: MAFFT --add, EPA-ng, gappa, iTOL v6.
% Type XI-like tree (Fig. 5): IQ-TREE2, LG+I+G4, 1,000 UFBoot, from
%   ``an alignment of the conserved RT domain'' --- span not specified, and not
%   stated to be the same 203 columns.
% Confidence tiers: 934 high (72.6%), 25 borderline (1.9%), 327 low (25.4%)
%   of 1,286 candidates. Criterion (ii), ``near-complete RT1--RT7 architecture'',
%   is not operationalised anywhere in the Methods: no threshold, no definition
%   of near-complete.
% Type XI-like: n = 232, 72% low-confidence; 207/232 (89.2%) lack the
%   C-terminal protease; 170/232 in Bacteroidota.
% Fig. 1 legend: outlying EPA-ng placements attributed to ``uncertainty in the
%   placement of highly divergent RT core sequences''.
% Minor internal inconsistencies in candidate counts: TXI_C2like 120 (Methods)
%   vs 126 (Results); TXI_noncan_h 28 vs 29; msr/msd hits 37 (text) vs 39
%   (Fig. 6B legend). Cite one consistently.


% =====================================================================
% =====================================================================
%  SECOND ROUND: additions following Supplementary Table S1
%  of \cite{toro2026} (HMM calibration thresholds, 29 profiles)
% =====================================================================
% =====================================================================


% ---------------------------------------------------------------------
% [EDIT] Footnote on supplementary material, \subsection{Origin...}
% Current text: ``...since a definition could reside in supplementary
% material.''  Extend the same convention to \cite{toro2026} with a
% checked statement.
% ---------------------------------------------------------------------

% Suggested footnote text for the \cite{toro2026} discussion:
%
% The same convention is applied here to~\cite{toro2026}: the statement is
% that no definition of the RT1--RT7 span appears in the published main text.
% Supplementary Table~S1 of that work reports the calibration of the
% type-specific profiles (self-hit minima, cross-hit maxima, per-group
% thresholds and separation status) and contains no element coordinates and
% no reference to RT0.


% =====================================================================
% [ADD] New paragraph for \subsection{Element-by-element}, following the
% RT0 discussion. This is the point at which the span enters the paper's
% own inferential chain.
% =====================================================================

The span is not merely a preprocessing choice in~\cite{toro2026}; it is named
as an input to the classification. Candidates are assigned to confidence tiers
on three stated criteria: HMM detection strength relative to the group
threshold, ``RT domain completeness, based on conservation of the canonical
RT1--RT7 motif core'', and consistency of phylogenetic placement. The second
criterion is the span. It is not operationalised anywhere in the published
methods---no threshold is given for completeness, no definition of
``near-complete'' is supplied, and no rule is stated for combining the three
criteria into a tier. The tiers in turn supply the ``confidence support''
column of the novelty evidence matrix, which contributes to the composite
score by which TXI\_C2like and TXI\_noncan\_h are selected as candidate novel
lineages. An element boundary that has never been derived therefore enters,
under a criterion that is never defined, the chain of inference leading to the
paper's principal claim.

Two features of Supplementary Table~S1 bear on whether the contribution of
that criterion can be recovered after the fact. The first is that the
supplementary legend defines the tiers by score alone: high-confidence as a
score at or above the group threshold, low-confidence as a score of at least
160 bits but below it, and borderline as a score falling within the overlap
between the self-hit and cross-hit distributions. Under that definition RT
domain completeness plays no part in the tier at all, which is not what the
main text states. The two descriptions cannot both be operative, and nothing
published distinguishes them. The second is that the borderline assignments do
not follow the supplementary definition either. All 25 borderline candidates
reported in the study fall in groups whose profiles are labelled as showing
clean separation (type~I-B2, $n = 12$; type~II-A2 and type~X, $n = 3$ each;
type~III-A4 and type~I-C3, $n = 2$ each; type~V, type~VII-A1 and type~VII-A2,
$n = 1$ each), and none falls in any of the groups labelled as overlapping.%
\todo{REG-07. Highest-risk claim in this section and the load-bearing one.
Derived by joining the Separation column of Supp.\ Table~S1 to the
borderline column of main Table~1. A single misread separation label or
borderline count flips it. Re-derive from the source supplementary file
before submission; do not cite from the extracted text.}
Since an overlap zone exists by construction only where the cross-hit maximum
exceeds the self-hit minimum, the observed assignment is the inverse of the
stated rule. Whatever produced the tiers, it is not the procedure described.
The consequence for the present chapter is narrow but firm: the weight carried
by the RT1--RT7 span in the classification of~\cite{toro2026} cannot be
determined from the published record, in either direction.


% =====================================================================
% [ADD] Footnote material --- criticisms of the classifier rather than of
% the partition. Recommended as a single footnote, not as body text, to
% keep the chapter on its subject.
% =====================================================================

% \todo{REG-01} Supplementary Table~S1 of~\cite{toro2026} reports 29 type-specific profiles,
% whereas the main text refers to a panel of 28 in each of the four places it
% is mentioned; the discrepancy is type~VIII, which is present in the panel,
% recovers no candidates, and does not appear in the summary table of results.
% \todo{REG-08}
% By the definition given in the same table---an overlap zone wherever the
% cross-hit maximum exceeds the self-hit minimum---17 of the 29 profiles fall
% in the overlap zone and 12 show clean separation; the totals row of the table
% states 16 and 13.  \todo{REG-02, REG-03} Each individual row is consistent with its own values. For
% the 17 overlapping profiles the classification threshold is the midpoint
% between two distributions that are not separated, and the main text reports
% only the range of thresholds (177.2--317.5 bits) rather than their
% separability. Type~XI is among the overlapping groups, and is the group from
% which both candidate novel lineages are drawn. Type~VI is the extreme case:
% its self-hit minimum, 162.9 bits, lies barely above the 160-bit general
% screening cutoff, 91 of its reference sequences fall below that cutoff, and
% its profile is in the overlap zone---yet 97\% of the type~VI candidates
% recovered from SPIRE are classified as high-confidence, which is what
% survivorship under a stringent screen would predict. Finally, the 195
% reference sequences excluded by the 160-bit screen decompose across seven
% groups (type~VI 91, type~I-B2 36, type~I-B1 34, type~IV 20, type~V 10,
% type~II-A1 2, type~VIII 2), of which the main text names four;  \todo{REG-05, REG-06} and no
% per-group reference totals are published, so the fraction of each group lost
% to the screen cannot be computed. That last omission is structurally the
% same as the one described in Section~\ref{sec:rt-partition-origin} for the
% 200-amino-acid threshold: a filter reported against an object whose extent
% is not stated.


% =====================================================================
% [ADD] Verified figures, second round. All recomputed from
% Supplementary Table S1 and main Table 1 of \cite{toro2026}.
% Values below were checked against the table's own definitions;
% confirm against the source supplementary file before citing.
% =====================================================================
% Panel size: 29 profiles (supp.) vs 28 (main text, four occurrences).
% Separation, recomputed from pos_min / neg_max: Overlap 17, Clean 12.
%   Printed totals row: Overlap 16, Clean 13. Every row label is itself correct.
% Threshold formula (pos_min + neg_max)/2 reproduces every printed cell.
%   clade11-A cell reads 317,4; range row and main text give 317.5 (=317.45).
%   Supplementary table uses decimal commas; main text uses points.
% Borderline candidates: 25 total, 25 in Clean groups, 0 in Overlap groups.
% Groups with zero recovered candidates: typeVIII (in panel, absent from Table 1).
% Refs below the 160-bit screen: 91 + 36 + 34 + 20 + 10 + 2 + 2 = 195 (10.1%),
%   across typeVI, typeI-B2, typeI-B1, typeIV, typeV, typeII-A1, typeVIII.
% No per-group reference denominators are published.
% clade2_Ec107like is annotated "new HMM in pipeline v2"; no pipeline
%   versioning is described in the main text.
% Threshold height does NOT explain confidence composition:
%   Pearson r(threshold, % high-confidence) = -0.28 across the 21 groups with
%   >= 10 candidates. Do not assert that the tier is an artefact of threshold
%   placement; it is not supported.


% =====================================================================
% =====================================================================
%  VERIFICATION REGISTER
%
%  Every quantity below was RECOMPUTED, not read directly off the page.
%  Source for all of it: text extracted from the main preprint PDF and
%  from Supplementary Table S1 (10.64898/2026.05.14.725207, v1,
%  14 May 2026).
%
%  GLOBAL CAVEAT: Supplementary Table S1 uses the comma as decimal
%  separator while the main text uses the point. Comma-decimal tables are
%  a known failure mode for PDF text extraction. Treat EVERY numeric cell
%  taken from S1 as provisional until re-derived from the source
%  supplementary file (ideally the XLSX/CSV, not the rendered PDF).
%
%  Anything NOT listed here is either quoted verbatim with a line number
%  (low risk) or is interpretation rather than measurement (defend it,
%  do not verify it).
% =====================================================================
%
% REG-01  Panel size = 29 profiles.
%         Method: counted label rows in S1; cross-checked against the
%         table title and the TOTALS row, both of which state 29.
%         Main text states 28 in four separate places.
%         Fails if: a row was lost in extraction.
%         Status: [ ] re-derived
%
% REG-02  Separation recount: Overlap 17, Clean 12.
%         Method: applied the table's own definition (overlap iff
%         neg_max > pos_min) to all 29 rows.
%         Printed TOTALS row states Overlap 16, Clean 13.
%         Every individual row LABEL agreed with its own values; only the
%         totals row disagrees.
%         Fails if: any single pos_min or neg_max cell misread.
%         Status: [ ] re-derived
%
% REG-03  Threshold formula reproduces all 29 printed cells.
%         Method: (pos_min + neg_max)/2 per row, compared to the printed
%         threshold column. This is the strongest internal consistency
%         check available, and it is what makes REG-02 credible: if the
%         extracted pos_min/neg_max values were corrupted, the formula
%         would not reproduce the printed thresholds.
%         Status: [ ] re-derived
%
% REG-04  clade11-A midpoint = 317.45.
%         The S1 cell reads 317,4; the S1 range row and the main text
%         both give 317.5. Rounding convention only. Do not build an
%         argument on this; mention in passing or omit.
%         Status: [ ] re-derived
%
% REG-05  Column G sums to 195: 91 + 36 + 34 + 20 + 10 + 2 + 2.
%         Reconciles exactly with the main text's stated 195 (10.1% of
%         1,928), which is independent corroboration that the seven G
%         cells were extracted correctly.
%         Status: [ ] re-derived
%
% REG-06  Seven groups contribute to the 195; the main text names four
%         (type VI, I-B2, I-B1, IV = 181 of 195), omitting type V (10),
%         type II-A1 (2) and type VIII (2).
%         Method: set difference between S1 column G and the main text
%         enumeration at l. 129-130.
%         Status: [ ] re-derived
%
% REG-07  *** HIGHEST RISK. LOAD-BEARING. ***
%         All 25 borderline candidates fall in Clean-separation groups;
%         zero fall in Overlap groups.
%         Method: join S1 Separation column to the borderline column of
%         main Table 1. Contributing groups: type I-B2 (12), type II-A2
%         (3), type X (3), type III-A4 (2), type I-C3 (2), type V (1),
%         type VII-A1 (1), type VII-A2 (1) = 25, matching the stated
%         total of 25 borderline candidates.
%         This is the inverse of the S1 legend, which defines borderline
%         as a score within the overlap zone -- a zone that exists by
%         construction only in Overlap groups.
%         Fails if: ONE separation label or ONE borderline count is
%         misread. There is no internal cross-check on this one.
%         RE-DERIVE THIS FIRST.
%         Status: [ ] re-derived
%
% REG-08  type VIII: present in the 29-profile panel, recovers zero
%         candidates, absent from main Table 1 (28 rows).
%         Method: set difference, S1 labels vs main Table 1 labels.
%         Status: [ ] re-derived
%
% REG-09  Pearson r(threshold, % high-confidence) = -0.275, over the 21
%         groups with >= 10 candidates.
%         NEGATIVE RESULT. Threshold height does NOT explain the
%         confidence composition. Recorded so the claim is not made by
%         accident; it is not supported.
%         Status: [ ] re-derived
%
% REG-10  Mean % high-confidence: Overlap groups 79.4 (n=16),
%         Clean groups 56.6 (n=5).
%         Direction is opposite to what the calibration design implies.
%         *** DO NOT PUT THIS IN THE THESIS AS IT STANDS. *** n=5 on the
%         Clean side, dominated by small groups (type I-C3 n=14, type X
%         n=16, type VII-A2 n=4). Suggestive only. If you want to use it,
%         weight by group size and report a confidence interval, or drop it.
%         Status: [ ] re-derived
%
% ---------------------------------------------------------------------
%  RECONCILIATION CHECKS THAT PASSED (no action needed, recorded for the
%  methods appendix if you want to show the audit was done)
% ---------------------------------------------------------------------
% Main Table 1 row sums: N = 1,286; high = 934; borderline = 25;
%   low = 327. All four match the stated totals.
% Mestre reference set: 1,912 + 16 validated = 1,928. Consistent across
%   \cite{mestre2020} and \cite{toro2026}.
%
% ---------------------------------------------------------------------
%  KNOWN INTERNAL INCONSISTENCIES IN \cite{toro2026} (read off the page,
%  not computed -- but pick one value each and cite consistently)
% ---------------------------------------------------------------------
% TXI_C2like:      120 loci (Methods, de novo ncRNA) vs 126 (Results, Fig. 4B)
% TXI_noncan_h:     28 loci (Methods, Abstract)      vs  29 (Results, curated)
% TXI_C2like ncRNA: 37 loci (Results text)           vs  39 (Fig. 6B legend)
% Panel size:       28 profiles (main text, x4)      vs  29 (Supp. Table S1)
% clade11-A thr:    317,4 (S1 cell)                  vs 317.5 (S1 range row,
%                                                        and main text)
%
% ---------------------------------------------------------------------
%  NOT VERIFIABLE QUANTITIES -- these are arguments, defend them, do not
%  "check" them
% ---------------------------------------------------------------------
% - The two readings of the "203 positions ... by removing columns
%   containing more than 50% gaps" sentence.
% - That no reason is given for the removal of RT0 (an absence claim over
%   the published main text; scope it explicitly, as you already do for
%   \cite{toro2019,mestre2020}).
% - The survivorship reading of type VI's 97% high-confidence fraction.
% - That the main-text and supplementary tier definitions are incompatible.
% - Simon & Zimmerly's 157 alignable retron characters vs the 203-column
%   block. DELIBERATELY NOT A COMPARISON: different quantities (cross-class
%   alignability to a group II reference vs within-retron gap occupancy).
%   State the gap; do not call it a contradiction.
%
% ---------------------------------------------------------------------
%  STILL TO CHECK IN THE SUPPLEMENT
% ---------------------------------------------------------------------
% [ ] Supp. Figure S1 (workflow diagram). Last plausible home for an
%     RT1-RT7 coordinate definition. If it shows the detection HMM and the
%     phylogeny alignment as separate branches, that confirms the
%     two-coordinate-system point directly rather than by inference.
% [ ] Supp. Table S2 (per-candidate table). Does it carry per-sequence
%     domain coordinates or a completeness score? If so, criterion (ii)
%     may be recoverable after all and the REG-07 argument needs softening.
% [ ] Supp. Table S3 (novelty scoring). How the five criteria combine into
%     the composite score.
% [ ] Per-group reference sequence totals. Without denominators the
%     fraction of each group lost to the 160-bit screen cannot be computed.
% =====================================================================
